% !TEX root =  master.tex
\chapter{Motivation und Zielsetzung}

\nocite{*}

\section{Motivation}

//TODO: Warum RNNs? Warum ist das Thema relevant? Limitierung der normalen Feedforward-Netze bei sequenziellen Daten, z.B. Zeitreihen

Trotz ihrer theoretischen Stärke leiden einfache \acp{RNN} in der Praxis unter dem Problem verschwindender und
explodierender Gradienten, sobald lange Abhängigkeitsketten erlernt werden müssen. Die Entwicklung von
\textit{Long Short-Term Memory}-Einheiten (LSTM) und \textit{Gated Recurrent Units} (GRU) hat dieses Problem
weitgehend gelöst und den Einsatzbereich rekurrenter Architekturen erheblich erweitert.

//TODO: Praxisbeispiele nennen, z.B. Sprachmodellierung, Zeitreihenanalyse, etc.

\section{Zielsetzung}

Das übergeordnete Ziel dieser Projektarbeit ist es, \acp{RNN} als Klasse sequenzieller Lernmodelle theoretisch
fundiert aufzuarbeiten und ihre praktischen Eigenschaften anhand von konkreten Beispielen zu veranschaulichen.

Hieraus leiten sich die folgenden Teilziele ab:

\begin{enumerate}
    \item \textbf{Theoretische Grundlagen:} Darstellung des mathematischen Modells einfacher rekurrenter Netze,
    einschließlich der Vorwärtsberechnung und der zeitlich entfalteten Rückwärtspropagation
    (\textit{Backpropagation Through Time}, BPTT).

    \item \textbf{Gradientenproblematik:} Analyse der Ursachen verschwindender und explodierender Gradienten
    sowie Diskussion der wichtigsten Gegenmaßnahmen.

    \item \textbf{Erweiterte Architekturen:} Vorstellung der GRU-Architekturen als Lösungsansatz
    für das Kurzeitgedächtnis-Problem einfacher \acp{RNN}.

    \item \textbf{Anwendungsbeispiele:} Prototypische Implementierung und Evaluation eines
    \ac{RNN}-basierten Modells an einem repräsentativen Datensatz zum Thema Aktienpreisvorhersage, um die erarbeiteten
    Konzepte zu veranschaulichen.

    \item \textbf{Kritische Einordnung:} Reflexion der Grenzen und Stärken rekurrenter Netze im Vergleich
    zu modernen Alternativen wie \textit{Transformer}-Architekturen.
\end{enumerate}

